\subsection{Frontend}
\label{sec:frontend}

Frontend to aplikacja jednostronicowa (SPA) napisana w React~19 i TypeScripcie
\cite{react}, budowana narzędziem Vite \cite{vite}. W Dockerze gotowe pliki
serwuje nginx (podrozdział~\ref{sec:docker}). Aplikacja komunikuje się z API
pod tym samym adresem (\texttt{/api}), więc nie wymaga konfiguracji CORS.

\subsubsection{Struktura aplikacji}

\begin{itemize}
  \item \texttt{api/}: klient HTTP (dodaje token, zamienia odpowiedzi błędów na
        wyjątek \texttt{ApiError} z kodem i błędami pól) oraz funkcje wywołujące
        poszczególne endpointy;
  \item \texttt{auth/}: stan zalogowania (kontekst Reacta), zapis sesji
        i strażnicy tras;
  \item \texttt{pages/}: ekrany aplikacji; \texttt{layout/}: układ stron;
        \texttt{components/}: wspólne elementy (pole formularza, komunikat);
  \item \texttt{i18n/texts.ts}: wszystkie teksty interfejsu po polsku w jednym
        module, więc komponenty nie zawierają napisów na stałe.
\end{itemize}

Nawigacją zarządza React Router, a adresy są po polsku (\texttt{/logowanie},
\texttt{/rejestracja}, \texttt{/profil}). Dane z serwera pobiera i przechowuje
w pamięci podręcznej biblioteka TanStack Query. Błędy klienta (4xx) nie są
ponawiane, bo ponowienie dałoby tę samą odpowiedź, a błędy sieci i serwera
są ponawiane jeden raz.

\subsubsection{Logowanie i sesja}

Po zalogowaniu token JWT, dane użytkownika i czas wygaśnięcia tokenu są
zapisywane w \texttt{localStorage}, dzięki czemu sesja przetrwa odświeżenie
strony i ponowne uruchomienie aplikacji zainstalowanej jako PWA. Przy starcie
zapisana sesja jest sprawdzana zapytaniem \texttt{GET /api/auth/me}. Każda
odpowiedź 401 na żądanie z tokenem (token wygasł albo konto usunięto) wylogowuje
użytkownika z komunikatem ,,Sesja wygasła. Zaloguj się ponownie.''. Aplikacja
wylogowuje też sama w chwili wygaśnięcia tokenu. Przechowywanie tokenu
w \texttt{localStorage} oznacza, że skrypt wstrzyknięty na stronę (atak XSS)
mógłby go odczytać. Ryzyko ogranicza React, który domyślnie zabezpiecza
wyświetlane dane, brak zewnętrznych skryptów i krótki czas życia tokenu
(1~godzina). Alternatywą byłoby ciasteczko \texttt{HttpOnly}, które wymaga
jednak dodatkowej ochrony przed atakami CSRF.

Strażnik \texttt{RequireAuth} przekierowuje niezalogowanego użytkownika do
ekranu logowania i zapamiętuje adres, pod który chciał wejść, a po zalogowaniu
kieruje go z powrotem pod ten adres. Strażnik \texttt{GuestOnly} nie pokazuje
ekranów logowania i rejestracji osobie już zalogowanej.

\subsubsection{Formularze}

Formularze obsługuje biblioteka react-hook-form ze schematami Zod
z pakietu \texttt{shared}. Te same reguły i te same polskie komunikaty
obowiązują więc w przeglądarce i w API. Niepoprawny formularz nie jest nawet
wysyłany, a błędy pojawiają się przy polach. Błędy, które zna tylko serwer
(np. zajęty login, odpowiedź 409), są przypisywane do właściwych pól na
podstawie \texttt{details} z odpowiedzi API. Po rejestracji aplikacja
automatycznie loguje użytkownika. Pola formularza są opisane etykietami,
a pole z błędem ma atrybut \texttt{aria-invalid} i powiązany komunikat
(\texttt{aria-describedby}), więc błąd odczytują też czytniki ekranu.

\subsubsection{Układ dla urządzeń mobilnych}

Style są pisane w podejściu \emph{mobile first}: podstawowe reguły dotyczą
telefonów, a szersze ekrany dostają poprawki w media query. Na telefonie
nawigacja główna jest paskiem przyklejonym do dolnej krawędzi ekranu, w zasięgu
kciuka. Wszystkie elementy dotykowe (przyciski, pola, odnośniki nawigacji) mają
co najmniej 48~pikseli wysokości, więcej niż zalecane minimum 44~pikseli.
Układ uwzględnia wycięcia ekranu (\texttt{env(safe-area-inset-*)}). Na ekranie
o szerokości od 768~pikseli nawigacja przechodzi pod górny pasek.

\douzupelnienia{lista i szczegóły dokumentów, profil, obsługa aparatu,
kadrowanie i kompresja (rozmiar przed i po), PWA, zrzuty ekranu z telefonu;
Etapy 8--10}
